
\exercise
Prove that if $S, T \in \L(V\1, W)$, then $\abs[\Big]{\, \norm{S} - \norm{T} \,} \le \norm{S-T}$. \label{5reverse}
\exercisecomment{The inequality above is called the \emph{reverse triangle inequality}.\index{reverse triangle inequality}}


\exercise
Suppose that $T \in \L(V)$ is self-adjoint or that $\F{}= \C{}$ and $T \in \L(V)$ is normal. Prove that
\[
\norm{T} = \max \br[\big]{ \abs{\la} : \la \text{ is an eigenvalue of } T}.
\]
\exercisedisplayend


\exercise
Suppose $T \in \L(V\1, W)$ and $v \in V$. Prove that\label{7eq}
\[
\norm{Tv} = \norm{T} \, \norm{v} \iff T^*T v = \norm{T}^2 v.
\]
\exercisedisplayend


\exercise
Suppose $T \in \L(V\1, W)$, $v \in V\1$, and $\norm{Tv} = \norm{T} \, \norm{v}$. Prove that if $u \in V$ and $\ip{u, v} = 0$, then
$\ip{Tu, Tv} = 0$.


\exercise
Suppose $U$ is a finite-dimensional inner product space, $T \in \L(V\1, U)$, and $S \in \L(U, W)$. Prove that \label{2prod}
\[
\norm{ST} \le \norm{S} \, \norm{T}.
\]
\exercisedisplayend


\exercise
Prove or give a counterexample: If $S, T \in \L(V)$, then $\norm{ST} = \norm{TS}$.


\exercise
Show that defining $d(S, T) = \norm{S-T}$ for $S, T \in \L(V\1, W)$ makes $d$ a metric on $\L(V\1, W)$.\label{7met}
\exercisecomment{This exercise is intended for readers who are familiar with metric spaces.}


\exercise
\begin{SUBtheorem}
\item
Prove that if $T \in \L(V)$ and $\norm{I - T} < 1$, then $T$ is invertible.
\item
Suppose that $S \in \L(V)$ is invertible. Prove that if $T \in \L(V)$ and $\norm{S - T} < 1/\norm[\big]{S^{-1}}$, then $T$ is invertible.
\end{SUBtheorem}
\exercisecomment{This exercise shows that the set of invertible operators in $\L(V)$ is an open subset of $\L(V)$, using the metric defined in Exercise \ref{7met}.}


\begin{comment}
\exercise
Suppose $S \in \L(V)$ is invertible. Prove that there exists $\delta > 0$ such that $T$ is invertible for every $T \in \L(V)$ with
$\norm{S-T} < \delta$.

\end{comment}

\exercise
Suppose $T \in \L(V)$. Prove that for every $\epsilon > 0$, there exists an invertible operator $S \in \L(V)$ such that
$0 < \norm{T - S} < \epsilon$. \label{7invert}


\exercise
Suppose $\dim V > 1$ and $T \in \L(V)$ is not invertible. Prove that for every $\epsilon > 0$, there exists
$S \in \L(V)$ such that $0 < \norm{T - S} < \epsilon$ and $S$ is not invertible.


\exercise
Suppose $\F{} = \C{}$ and $T \in \L(V)$. Prove that for every $\epsilon > 0$ there exists a diagonalizable operator $S \in \L(V)$ such that $0 < \norm{T - S} < \epsilon$.\index{diagonalizable}


\exercise
Suppose $T \in \L(V)$ is a positive operator. Show that $\norm[\Big]{\sqrt{T}\,} = \sqrt{\norm{T}}$.


\exercise
Suppose $S, T \in \L(V)$ are positive operators. Show that
\[
\norm{S-T} \le \max\br[\big]{ \norm{S}, \norm{T} } \le \norm{S+T}.
\]
\exercisedisplayend


\exercise
Suppose $U$ and $W$ are subspaces of $V$ such that $\norm{P_U - P_W} < 1$. Prove that $\dim U = \dim W$.


\pagebreak[3]

\exercise
Define $T \in \L\paren[\big]{\F{3}}$ by
\[
T(z_1, z_2, z_3) = (z_3, 2z_1, 3z_2).
\]
Find (explicitly) a unitary operator $S \in \L\paren[\big]{\F{3}}$ such that $T = S \sqrt{T^*T}$.


\exercise
Suppose $S \in \L(V)$ is a positive invertible operator. Prove that there exists $\delta > 0$ such that $T$ is a positive operator for every self-adjoint operator $T \in \L(V)$ with $\norm{S-T} < \delta$.


\exercise
Prove that if $u \in V$ and $\phi_u$ is the linear functional on $V$ defined by the equation $\phi_u(v) = \ip{v, u}$, then
$\norm{\phi_u} = \norm{u}$.
\exercisecomment{Here we are thinking of the scalar field\/ $\F{}$ as an inner product space with\/ $\ip{\al,\beta} = \al \overline{\beta}$ for all\/ $\al, \beta \in \F{}\3$. Thus\/ $\norm{\phi_u}$ means the norm of\/ $\phi_u$ as a linear map from\/ $V$ to\/ $\F{}\3$.}


\exercise
Suppose $e_1, \ldots, e_n$ is an orthonormal basis of $V$ and $T \in \L(V\1,W)$.
\begin{subtheorem}
\item
Prove that
$\max\br{ \norm{Te_1}, \ldots, \norm{Te_n} }
\le \norm{T} \le \paren[\big]{ \norm{Te_1}^2 + \cdots + \norm{Te_n}^2 }^{1/2}\!$.
\item
Prove that $\norm{T} = \paren[\big]{ \norm{Te_1}^2 + \cdots + \norm{Te_n}^2 }^{1/2}$ if and only if $\dim \ran T \le 1$.
\end{subtheorem}
\exercisecomment{Here\/ $e_1, \ldots, e_n$ is an arbitrary orthonormal basis of\/ $V\1$, not necessarily connected with a singular value decomposition of\/ $T\3$. If\/ $s_1, \ldots, s_n$ is the list of singular values of\/ $T$, then the right side of the inequality above equals\/
$\paren[\big]{{s_1}^{\2} + \cdots + {s_n}^{\2}}^{1/2}\!$, as was shown in Exercise \ref{7sums2}\uppar{a} in Section \ref{2SVD}.}


\exercise
Prove that if $T \in \L(V\1, W)$, then  $\norm[\big]{T^*T} = \norm{T}^2\1$.\index{C*@$C^*$-algebras}\label{8Cstar}
\exercisecomment{This formula for\/ $\norm[\big]{T^*T}$ leads to the important subject of\/ $C^*$-algebras.}


\exercise
Suppose $T \in \L(V)$ is normal. Prove that $\norm[\big]{T^k} = \norm{T}^k$ for every positive integer $k$.\label{7normTk}


\exercise
Suppose $\dim V > 1$ and $\dim W > 1$. Prove that the norm on $\L(V\1, W)$ does not come from an inner product. In other words, prove that there does not exist an inner product on $\L(V\1, W)$ such that \label{6notinner}
\[
\max \br[\big]{ \norm{Tv} : v \in V \text{ and } \norm{v} \le 1} = {\textstyle \sqrt{\ip{T, T}}}
\]
for all $T \in \L(V\1, W)$.


\exercise
Suppose $T \in \L(V\1, W)$. Let $n = \dim V$ and let $s_1 \ge \cdots \ge s_n$ denote the singular values of $T\3$. Prove that if $1 \le k \le n$, then \label{7approx}
\[
\min \br[\big]{ \norm{T|_U} : U \text { is a subspace of } V \text{ with } \dim U = k} = s_{n-k+1}.
\]
\exercisedisplayend


\exercise
Suppose $T \in \L(V\1, W)$. Show that $T$ is uniformly continuous with respect to the metrics on $V$ and $W$ that arise from the norms on those spaces (see Exercise \ref{4metric} in Section \ref{orthobases}).


\exercise
Suppose $T \in \L(V)$ is invertible. Prove that
\vspace{-3bp}
\[
\norm[\big]{ T^{-1} } = \norm{T}^{-1} \iff \frac{T}{\norm{T}} \text{ is a unitary operator}.
\]
\vspace{-16bp}


\pagebreak[3]

\exercise
Fix $u, x \in V$ with $u \ne 0$. Define $T \in \L(V)$ by
$
Tv = \ip{v, u}x
$
for every $v \in V\1$.  Prove that
\[
\sqrt{T^* T}v = \frac{ \|x\| }{ \|u\| } \ip{v,u} u
\]
for every $v \in V\1$.


\exercise
Suppose $T \in \L(V)$.  Prove that $T$ is invertible if and only if there exists a
unique unitary operator $S \in \L(V)$ such that $T = S \sqrt{T^* T}$.


\exercise
Suppose $T \in \L(V)$ and $s_1, \ldots, s_n$ are the singular values of $T\3$. Let $e_1, \ldots, e_n$ and $f_1, \ldots, f_n$ be orthonormal bases of $V$ such that \label{7unitary}
\[
Tv = s_1\ip{v, e_1} f_1 + \cdots + s_n \ip{v, e_n} f_n
\]
for all $v \in V\1$. Define $S \in \L(V)$ by
\[
Sv = \ip{v, e_1} f_1 + \cdots + \ip{v, e_n} f_n.
\]
\begin{subtheorem}*
\item
Show that $S$ is unitary and $\norm{T-S} = \max \br{ \abs{s_1 -1}, \ldots, \abs{s_n - 1} }$.
\item
Show that if $E \in \L(V)$ is unitary, then $\norm{T-E} \ge \norm{T-S}$.
\end{subtheorem}
\exercisecomment{This exercise finds a unitary operator $S$ that is as close as possible \uppar{among the unitary operators} to a given operator $T\3$.}


\exercise
Suppose $T \in \L(V)$. Prove that there exists a unitary operator $S \in \L(V)$ such that
$
T = \sqrt{TT^*}\, S$.


\exercise
Suppose $T \in \L(V)$.
\begin{subtheorem}
\item
Use the polar decomposition to show that there exists a unitary operator $S \in \L(V)$ such that
$TT^* = ST^*TS^*\1$.
\item
Show how (a) implies that $T$ and $T^*$ have the same singular values.
\end{subtheorem}
\exercisesubtheoremend


\exercise
Suppose $T \in \L(V)$, $S \in \L(V)$ is a unitary operator,
and $R \in \L(V)$ is a  positive operator such that $T = SR$.  Prove that
$R = \sqrt{T^* T}$.
\exercisecomment{This exercise shows that if we write\/ $T$ as the product of a unitary operator and a positive operator
\uppar{as in the polar decomposition \ref{244}\hspace{.6bp}}, then the positive operator equals\/
$\sqrt{T^*T}$.}


\exercise
Suppose $\F{} = \C{}$ and $T \in \L(V)$ is normal. Prove that there exists a unitary operator $S \in \L(V)$ such that $T = S \sqrt{T^*T}$ and such that $S$ and $\sqrt{T^*T}$ both have diagonal matrices with respect to the same orthonormal basis of $V\1$. \label{8PolarNor}


\exercise
Suppose that $T \in \L(V\1,W)$ and $T \ne 0$. Let $s_1, \ldots, s_m$ denote the positive singular values of $T\3$. Show that there exists an orthonormal basis $e_1, \ldots, e_m$ of $(\ker T)^\perp$ such that
\[
T\paren[\bigg]{ E\paren[\Big]{ \frac{e_1}{s_1}, \ldots, \frac{e_m}{s_m} } }
\]
equals the ball in $\ran T$ of radius $1$ centered at $0$.


